\section{Corrected volume holonomy, affine density, and deep-prime escape}
\label{sec:holonomy-density-depth}

This section records the next continuation of the four active routes.  The
results are deliberately asymmetric.  On the positive side we construct an
inhabited valuation-ball model for the local volume law, prove an affine
squarefree-bulk theorem, and amplify the arithmetic consequences of a
hypothetical Pell or Danilov survivor.  On the negative side we give complete
counterexamples to two scalar reconstruction mechanisms.  Neither
counterexample satisfies the stronger object-level hypotheses of the
corresponding IUT or arithmetic route.

\subsection{A corrected local volume model and its holonomy ledger}

Fix a rational prime $p$, normalize $|p|_p=p^{-1}$, and for $k\in\Z$ put
\[
             B_p(k)=\{x\in\Q_p:|x|_p\leq p^{-k}\}.
\]
These compact-open balls avoid the empty-set obstruction in
Theorem~\ref{thm:gpsp-lana-empty} while retaining the intended scaling law.

\begin{proposition}[Valuation-ball transport]
\label{prop:hdd-ball-transport}
For every $k\in\Z$,
\[
       (x\mapsto px)^{-1}B_p(k)=B_p(k-1).
\]
Moreover, $B_p(k)\subseteq B_p(m)$ if and only if $m\leq k$.  If
$\operatorname{vol}(B_p(k))=p^{-k}$, then
\[
 \log\operatorname{vol}((x\mapsto px)^{-1}B_p(k))
   =\log\operatorname{vol}(B_p(k))+\log p.
\]
The same additive shift survives every normalized finite weighted packet and
every positive-length procession average.
\end{proposition}

\begin{proof}
The first assertion is the chain
\[
 |px|_p\leq p^{-k}
 \iff p^{-1}|x|_p\leq p^{-k}
 \iff |x|_p\leq p^{-(k-1)}.
\]
The inclusion criterion follows because $p^{-k}\leq p^{-m}$ precisely when
$m\leq k$; necessity may also be tested at $p^k$.  The volume formula is
$\log p^{-(k-1)}=\log p^{-k}+\log p$.  Finally, a common additive shift
passes through a weighted sum whose weights total one, and through division
by a positive finite procession length.
\end{proof}

Let a logarithmic transport from a pointed object $x$ to a pointed object
$y$ with coordinate $L$ and shift $\delta$ mean
\[
                              L(y)=L(x)+\delta.
\]
Composition adds shifts and reversal negates them.

\begin{theorem}[Zero logarithmic holonomy]
\label{thm:hdd-zero-holonomy}
Every transport returning to the same pointed object has total shift zero.
Consequently an uncorrected loop whose edge shifts are nonnegative and whose
total shift is positive cannot be a same-pilot loop.  A correcting edge in a
closed loop must contribute the negative of the accumulated shift.
\end{theorem}

\begin{proof}
If $y=x$, then $L(x)=L(x)+\delta$, and cancellation gives $\delta=0$.
The remaining assertions are immediate applications to a sum of edge
shifts.
\end{proof}

For rational coefficients the scalar test can be sharpened place by place.

\begin{theorem}[Prime-local rational ledger]
\label{thm:hdd-prime-ledger}
Let $p_1,\ldots,p_r$ be distinct rational primes and $c_i\in\Q$.  If
\[
                         \sum_{i=1}^r c_i\log p_i=0,
\]
then $c_i=0$ for every $i$.  Hence a closed rational-prime logarithmic
transport must cancel separately at every rational place.
\end{theorem}

\begin{proof}
Clear a common denominator and write the resulting coefficients as
$m_i\in\Z$.  Exponentiation gives $\prod_i p_i^{m_i}=1$.  Moving negative
powers to the other side produces an equality of positive integers, and
unique factorization forces every $m_i$ to vanish.
\end{proof}

The coefficient field and the amount of retained local data are essential.
Arbitrary real coefficients are dependent, for example
$(\log3)\log2-(\log2)\log3=0$.  Positivity and normalization do not repair
the loss of labelled information.

\begin{proposition}[Positive normalized scalar collision]
\label{prop:hdd-positive-collision}
There are two different triples of strictly positive real weights, each
having sum one, whose weighted averages of $\log2,\log3,\log5$ are equal.
\end{proposition}

\begin{proof}
Put $a=\log2<b=\log3<c=\log5$ and define
\[
 W=\left(\frac{c-b}{2(c-a)},\frac12,
          \frac{b-a}{2(c-a)}\right),\qquad
 V=\left(\frac{c-b}{3(c-a)},\frac23,
          \frac{b-a}{3(c-a)}\right).
\]
All six entries are positive, and each row sums to one.  More generally, for
$n=2,3$,
\[
 \frac{c-b}{n(c-a)}a+\left(1-\frac1n\right)b
       +\frac{b-a}{n(c-a)}c=b,
\]
because $(c-b)a+(b-a)c=b(c-a)$.  The middle coordinates $1/2$ and $2/3$
show that the triples differ.
\end{proof}

Proposition~\ref{prop:hdd-positive-collision} is a full counterexample to
recovering labelled weights from one positive normalized scalar.  It does
not refute a construction which transports the labelled local objects, nor
the stronger same-pilot map equality discussed after
Proposition~\ref{prop:gpsp-pointed-hit}.

\subsection{Squarefree bulk and the coupled long-arm affine gate}

Return to the minimal radical-step shear in
Theorem~\ref{thm:gpsp-radical-shear}.  Thus $a+b=c$ is primitive,
$R=\rad(abc)$, and
\[
 U=1+Rh,\qquad V=1+R(h+ck),\qquad W=1+R(h+bk).
\]

\begin{proposition}[Support-closed gap recovery]
\label{prop:hdd-gap-recovery}
For $h,k>0$,
\[
 \gcd(V-U,W-U)=Rk,
 \qquad
 (a,b,c)=\frac1{Rk}(V-W,W-U,V-U).
\]
Conversely, let $1<U<W<V$, put $g=\gcd(V-U,W-U)$, and divide the
three consecutive gaps by $g$ to obtain $(a,b,c)$.  If
$\rad(abc)\mid g$ and $\rad(abc)\mid U-1$, then the triple is exactly a
minimal positive-parameter radical-step shear with $k=g/R$ and
$h=(U-1)/R$.
\end{proposition}

\begin{proof}
The forward identities follow from
$(V-W,W-U,V-U)=Rk(a,b,c)$ and $\gcd(b,c)=1$.  Conversely the definition of
$g$ makes the normalized gaps integral and primitive, while addition of the
first two gives the third.  The two divisibilities define integral $h,k$;
the strict inequality $U>1$ gives $h>0$, and substitution recovers the
three displayed affine forms.  The weaker condition $U\geq1$ is false:
$(U,W,V)=(1,3,5)$ gives $(a,b,c)=(1,1,2)$ and $R=g=2$, but $h=0$.
\end{proof}

Write $E(n)=n/\rad(n)$.  The inherited three-quarter exceptional inequality
in the upper-half canonical box is
\[
                         E(U)E(V)E(W)>\frac{Rc^{14}}{8192},
\]
while $U\leq c^6$ and $V,W\leq c^7$.  It follows pointwise that
\begin{equation}\label{eq:hdd-two-long-arms}
                         8192E(V)>Rc,\qquad8192E(W)>Rc.
\end{equation}
Thus both long arms must have large repeated-prime excess.

The same box nevertheless contains a uniformly positive squarefree bulk.
Put
\[
 M=\left\lfloor\frac{c^6}{4R}\right\rfloor,
 \qquad I_M=\{\lfloor M/2\rfloor+1,\ldots,M\},
 \qquad N=|I_M|.
\]

\begin{theorem}[Uniform squarefree-admissible bulk]
\label{thm:hdd-affine-squarefree-bulk}
Fix $\theta<5/2$.  For all sufficiently large primitive seeds satisfying
$R<c^\theta$, at least
\[
       \frac12\left(5-\frac{\pi^2}{2}\right)N^2
\]
pairs $(h,k)\in I_M^2$ obey $\gcd(U,k)=1$ and have $U,V,W$ all squarefree.
Every resulting $abc$ output has radical at least its height and is therefore
nonexceptional for every fixed exponent below one.
\end{theorem}

\begin{proof}
For a prime $p\nmid R$, failure of $\gcd(U,k)=1$ prescribes one residue class
for both $h$ and $k$ modulo $p$, and hence occurs for at most
$(N/p+1)^2$ pairs.  For a fixed $F\in\{U,V,W\}$, the condition $p^2\mid F$
prescribes one class of $h$ modulo $p^2$ after $k$ is fixed, and hence occurs
for at most $N^2/p^2+N$ pairs.  Primes dividing $R$ cause neither event.
Since $2\mid R$, the union bound and
\[
 \sum_{j\geq1}\frac1{(2j+1)^2}=\frac{\pi^2}{8}-1
\]
leave at least
\[
 \left(5-\frac{\pi^2}{2}\right)N^2
       -2N(1+\log M)-M-3Nc^{7/2}
\]
good pairs.  Here every possible square divisor is covered because
$U,V,W\leq c^7$.  The hypothesis $R<c^\theta$ gives
$N\gg c^{6-\theta}$, so the three error terms are smaller than half the
positive main term when $\theta<5/2$.

For a good pair, pairwise coprimality and avoidance of the seed support give
$\rad(aUbVcW)=RUVW$.  Also $U\geq c$ for all sufficiently large seeds in
this range, and therefore $RUVW\geq cW$, the output height.
\end{proof}

Theorem~\ref{thm:hdd-affine-squarefree-bulk} and
\eqref{eq:hdd-two-long-arms} point in opposite directions without
contradiction.  The desired exceptional lower count occupies a vanishing
fraction of the raw box, so a positive-density generic bulk cannot be
complemented into that lower bound.  A proof must create simultaneous
high-excess tails in both long arms or exploit the support-closed reverse
description in Proposition~\ref{prop:hdd-gap-recovery}.

\subsection{Prime-rank depth and second-order Pell coupling}

Let $u_0=0,u_1=1,u_{n+2}=6u_{n+1}-u_n$ and write
$(1+\sqrt2)^n=A_n+B_n\sqrt2$.  At an odd prime index $\ell$,
$u_\ell=A_\ell B_\ell$, and the inherited valuation dictionary identifies
every odd divisor of either channel with rank $\ell$.

\begin{theorem}[Pointwise channel escape]
\label{thm:hdd-channel-escape}
For every odd prime $\ell$, the $A$ channel has either a prime
$q\parallel A_\ell$ or an odd-exponent divisor $p^{e(p)}\mid A_\ell$ with
rank $\ell$ and $e(p)\geq3$.  The same holds for the $B$ channel except for
$B_7=13^2$.  Consequently each channel separately has simple divisors at all
but finitely many prime indices, or has infinitely many distinct prime-rank
depth-three primes.  If the global set of depth-three primes is finite, then
both simple divisors occur simultaneously at every sufficiently large prime
index.
\end{theorem}

\begin{proof}
An integer which is not a perfect power either has an exponent-one prime or
has an odd prime exponent at least three.  Apply this observation to the
perfect-power classifications of the two Pell coordinates.  Different prime
indices give different ranks, hence different primes in the second branch.
If only finitely many depth-three primes exist, delete their finitely many
ranks and the exceptional index seven; both channels must then take the
simple branch at every remaining prime index
\cite{Cohn1996Pell,Cohn2003Pell,Ljunggren1942}.
\end{proof}

Factoring the channels and expanding every signed local factor through its
quadratic term yields a necessary congruence one digit deeper than the
first-order ledger.  With the notation $K_A,K_B,C_A,C_B$ for the linear and
quadratic factor coefficients, respectively,
\begin{equation}\label{eq:hdd-second-order}
 K_A-2K_B+\ell(K_A^2-2K_B^2)
       +2\ell(C_A-2C_B)\equiv0\pmod{4\ell^2}.
\end{equation}
Quadratic reciprocity also gives the all-pairs identity
\[
 \prod_{q\mid A_\ell}\prod_{r\mid B_\ell}
       \left(\frac qr\right)^{v_q(A_\ell)v_r(B_\ell)}
       =\left(\frac2\ell\right).
\]
These ledgers constrain a full four-prime packet but do not yet force an
inconsistency.

Applying the repaired Fellini--Murty prime-order argument separately to
$1+\sqrt2$ and its negative produces an alternative: either infinitely many
rational balancing primes have depth at least three, or infinitely many
prime indices have a simple $A$-divisor and infinitely many prime indices
have a simple $B$-divisor.  The latter two index sets are not asserted to
intersect; the synchronized conclusion in
Theorem~\ref{thm:hdd-channel-escape} uses the stronger finite-depth-three
hypothesis instead~\cite{FelliniMurty2026}.

An exhaustive certificate checks every one of the $50{,}847{,}533$ odd
primes $q\leq10^9$.  Independent segmented-sieve/Lucas-doubling and
dense-sieve/quotient-ring implementations agree: the only depth-two primes
are
\[
                         13,\quad31,\quad1546463,
\]
and each has exact depth two; no depth-three prime occurs.  Arbitrary-precision
replay verifies the three residues modulo $q^3$.  Thus each depth-three prime
required by an actual opposite-channel packet exceeds $10^9$.  This is a
finite lower bound, not a nonexistence result.

\subsection{Divisor-pair amplification on the Danilov survivor}

Let $F_n$ be the Fibonacci sequence.  The Danilov factorization has the
form
\[
                    K_t=\frac{27}{25}F_NF_{N-5},
              \qquad N=10(3t+1),
\]
with $\gcd(F_N,F_{N-5})=5$.  The final recursive certificate supplies a
squarefree, $638$-prime, $4398$-digit modulus $Q_*$, coprime to $30$, such
that $Q_*\mid3t+1$ throughout the surviving progression.

\begin{theorem}[Wall--Sun--Sun divisor-pair amplification]
\label{thm:hdd-wss-amplification}
Let $Q>1$ be squarefree and coprime to $10$, let $R$ be a positive multiple
of $Q$, and put $N=10R$.
If $(27/25)F_NF_{N-5}$ is squarefull and $Q$ has $s$ distinct prime factors,
then there are at least $2^{s-1}$ distinct Wall--Sun--Sun primes.  They may
be indexed by the divisors $d\mid Q$ satisfying $d<Q/d$, and the indexed
prime $p_d$ obeys
\[
                z(p_d)=N/d,\qquad p_d\equiv1\pmod{N/d},
                \qquad p_d>N/\sqrt Q.
\]
For the final Danilov progression this supplies a distinguished subfamily of
at least $2^{637}$ such primes, each exceeding $10^{2199}$.
\end{theorem}

\begin{proof}
The involution $d\mapsto Q/d$ pairs the $2^s$ divisors of nonsquare $Q$, so
exactly $2^{s-1}$ satisfy $d<Q/d$.  Put $m_d=N/d$.  Carmichael--Yabuta gives
a primitive divisor $p_d$ of $F_{m_d}$~\cite{Yabuta2001}.  Since
$10\mid m_d$, its rank is
$m_d$ and the split-rank congruence gives $p_d\equiv1\pmod{m_d}$.  The
inequality $d<\sqrt Q$ implies $m_d>d$, whence $p_d$ divides neither $m_d$
nor $d$ and therefore does not divide $N=m_dd$.

Now $p_d\mid F_N$ and, because $p_d>5$, it does not divide $F_{N-5}$ or the
rational prefactor.  Squarefullness forces $v_{p_d}(F_N)\geq2$.  The exact
Lucas valuation formula and $p_d\nmid N$ transfer this valuation to
$F_{m_d}$, making $p_d$ a Wall--Sun--Sun prime~\cite{Sanna2016}.  Equality of two selected
primes would force equality of their ranks and hence equality of the
corresponding divisors, so they are distinct.  Finally
$p_d\geq N/d+1>N/\sqrt Q$.

For $Q_*$ take $s=638$.  Since $Q_*$ has $4398$ decimal digits and
$R\geq Q_*$, the size bound gives
$p_d>10\sqrt{Q_*}>10^{2199}$.
\end{proof}

The certified upper bound for the prime factors of $Q_*$ permits almost all
divisors, rather than only one representative from each complementary pair.

\begin{theorem}[Factor-bound amplification]
\label{thm:hdd-wss-factor-bound}
Under the hypotheses of
Theorem~\ref{thm:hdd-wss-amplification}, suppose every prime factor of $Q$
is at most $B$, and put
\[
 E_B(Q)=\#\{e:e\mid Q,\ 10e\leq B\}.
\]
Then squarefullness forces at least $2^s-E_B(Q)$ distinct Wall--Sun--Sun
primes.  At the verified endpoint $B=99{,}966{,}059$ and exact truncated
subset-product enumeration gives $E_B(Q_*)=622$.  Thus a survivor forces at
least $2^{638}-622$ distinct Wall--Sun--Sun primes.  At least one forced
prime is at least $41N+1>10^{4399}$.
\end{theorem}

\begin{proof}
Write $R=kQ$.  For every divisor $e\mid Q$ with $10e>B$, put
$d=Q/e$ and $m_e=N/d=10ke$.  A primitive divisor $p_e$ of $F_{m_e}$ is
congruent to one modulo $m_e$, so $p_e>m_e>B$.  It divides neither $Q$ nor
$m_e$, and hence does not divide $N=m_ed$.  The squarefullness and valuation
transfer argument from Theorem~\ref{thm:hdd-wss-amplification} makes it a
Wall--Sun--Sun prime.  Different $e$ give different ranks and hence different
primes.  Exactly $E_B(Q)$ of the $2^s$ divisors were omitted.

For $Q_*$ the certified factor list has maximum $B=99{,}966{,}059$; exact
enumeration below $B/10$ returns $622$ divisors.  The final large-prime claim
uses Hong's primitive-divisor theorem with $\kappa=40$: since
$\log N>10036$, it supplies a primitive divisor of $F_N$ outside
$jN\mathbin\pm1$ for $1\leq j\leq40$.  The split congruence leaves the plus
sign and coefficient at least $41$, and squarefullness transfers repeatedness
as before~\cite{Hong2025}.
\end{proof}

No known theorem excludes $2^{638}-622$ Wall--Sun--Sun primes at these sizes,
so Theorems~\ref{thm:hdd-wss-amplification} and
\ref{thm:hdd-wss-factor-bound} are strong, explicitly quantified conditional
obstructions,
not a contradiction.  The cyclotomic-derivative shortcut also fails: the
complete example
\[
 \Phi_{10}(-3)=121=11^2,\qquad \Phi_{10}'(-3)=-142,
 \qquad \gcd(121,142)=1
\]
shows that a squared value need not come from a multiple root.  Hensel
uniqueness does not prevent a fixed algebraic unit from landing on the unique
lift.  This counterexample closes only that derivative mechanism; the
Fibonacci and Danilov routes remain active.

\subsection{Formal and logical boundary}

The elementary parts of Propositions~\ref{prop:hdd-ball-transport} and
\ref{prop:hdd-positive-collision}, Theorems~\ref{thm:hdd-zero-holonomy} and
\ref{thm:hdd-prime-ledger}, the affine gap and excess identities, the Pell
second-order ledger, and the abstract Danilov divisor-pair transfer have
independent Lean implementations.  Across the four companion modules these
comprise $65$ theorem or lemma declarations, $18$ definitions, four
abbreviations, and five structures, for $92$ declarations; $58$ embedded
\texttt{\#print axioms} queries report only \texttt{propext},
\texttt{Classical.choice}, and \texttt{Quot.sound}.  The squarefree-bulk
estimate, the perfect-power classifications and Fellini--Murty argument, the
primitive-divisor and valuation theorems, Hong's theorem, and the large finite
certificates remain explicit paper, literature, or computation inputs and are
not introduced as Lean axioms.  The computation packages freeze source,
output, byte count, and SHA-256 manifests.

At this stage the four live endpoints are: construct an object-level
same-pilot transport satisfying the prime-local zero-holonomy ledger; produce
or exclude the simultaneous affine high-excess tails; rule out or construct
opposite-channel prime-rank depth-three Pell packets; and exclude or realize
the enormous Wall--Sun--Sun population forced by a Danilov survivor.  None is
closed by difficulty or by a finite no-hit, and no unconditional term of type
\texttt{ABCConjecture} or its negation is claimed.
